\documentclass[10pt,a4paper]{article}
\usepackage[T1]{fontenc}
\usepackage[margin=2cm]{geometry}
\usepackage{isabelle,isabellesym}
\usepackage{amsmath,amssymb}
\usepackage{pdfsetup}

\urlstyle{rm}
\isabellestyle{it}

\begin{document}

\title{Tails of Sums, Products and Random Sums of Regularly Varying Random Variables}
\author{Carles Mar\'in}
\maketitle

\begin{abstract}
Building on the AFP entries on regular variation and extreme value theory, we formalize the basic tail calculus of
heavy-tailed random variables~\cite{BinghamGoldieTeugels1987,EmbrechtsKluppelbergMikosch1997}. For independent $X$ and $Y$, if
$P(X>x)$ is regularly varying and $P(Y>x)/P(X>x)\to c$, then $P(X+Y>x)/P(X>x)\to 1+c$. Breiman's lemma~\cite{Breiman1965}: if $X\ge 0$ has a regularly varying
tail of index $-\alpha$, $Z\ge 0$ is independent of $X$ and $E\,Z^{\alpha+\varepsilon}<\infty$, then
$P(XZ>x)\sim E[Z^\alpha]\,P(X>x)$. We deduce closure of the Fr\'echet domain of attraction under these operations, the
big-jump principle $P(S_n>x)\sim P(\max_i X_i>x)\sim n\,P(X_1>x)$, and an explicit Pareto example showing that independence
cannot be dropped. For a random number $N$ of summands, independent of the sequence, with
$E\,N^{1+\alpha+\delta}<\infty$, we prove $P(S_N>x)\sim E[N]\,P(X_1>x)$, in particular for geometric $N$. Finally, for a
service time $B$ with regularly varying tail of index $-\alpha$, $\alpha>1$, we prove the regular variation of the integrated tail
and, \emph{assuming} the Pollaczek--Khinchine representation of the stationary waiting time $W$ as a geometric sum, the
classical asymptotics $P(W>x)\sim \frac{\rho}{1-\rho}\,\overline{B}_I(x)$~\cite{Pakes1975}. The queue itself is not constructed:
that representation is a hypothesis of the theorem.
\end{abstract}

\paragraph{Acknowledgement.} The formalization and the text of this document were prepared with the
assistance of generative AI systems, and checked by the author.

\tableofcontents

\parindent 0pt\parskip 0.5ex

\input{session}

\bibliographystyle{abbrv}
\bibliography{root}

\end{document}
